% Figure 9-A -- Component symbol plate: the cognitive inductor.
% (a) L1 four-loop coil with attention-flow terminals + parameter box
% (b) mutual-inductance pair sharing a dashed core, labelled M and k
% (c) flyback diode wired across a coil, reverse-blocking orientation
\begin{tikzpicture}[x=1cm,y=1cm,font=\sffamily\footnotesize,
  wire/.style={line width=0.7pt, draw=black},
  coilw/.style={line width=0.9pt, draw=casCedge},
  panlbl/.style={font=\sffamily\scriptsize\bfseries, text=casNedge, anchor=west},
  tny/.style={font=\sffamily\scriptsize}]

% =============== (a) the cognitive inductor =========================
\node[panlbl] at (0.35,1.00) {(a)\ \ COGNITIVE INDUCTOR L1};

\node[casterm, text width=15mm] (tin)  at (1.15,0) {attention\\flow in};
\node[casterm, text width=15mm] (tout) at (7.90,0) {attention\\flow out};

\draw[flow] (tin.east) -- (2.98,0);
\draw[flow] (6.02,0) -- (tout.west);

% four-loop coil
\draw[coilw] (3.05,0) -- (3.55,0);
\foreach \i in {0,1,2,3}{\draw[coilw] ({3.55+0.50*\i},0) arc(180:0:0.25);}
\draw[coilw] (5.55,0) -- (6.05,0);
\fill (3.05,0) circle (1.2pt);
\fill (6.05,0) circle (1.2pt);

\node[font=\sffamily\footnotesize\bfseries] at (4.55,0.66) {L1};

\node[nblock, text width=70mm, align=left, anchor=north,
      font=\sffamily\scriptsize] at (4.55,-0.58) {%
  inductance $L$ in h$\cdot\Omega_c$:\ \ $L = k_d \times H_{\mathrm{practice}}$\\
  opposes the \emph{rate}, not the flow:\ \ $V_{\mathrm{opp}} = L\,|di/dt|$\\
  depth coefficient $k_d$ = 0.5 / 1.0 / 2.0 / 4.0 for O/E/C/P};

% ------------- divider ---------------------------------------------
\draw[draw=casNedge, dash pattern=on 1.5pt off 1.5pt, line width=0.4pt]
  (0.3,-2.60) -- (8.8,-2.60);

% =============== (b) mutual inductance ==============================
\node[panlbl] at (0.35,-2.94) {(b)\ \ MUTUAL INDUCTANCE};

% left coil L1, bumps facing away from the shared core
\draw[coilw] (1.60,-4.95) -- (1.60,-4.80);
\foreach \i in {0,1,2,3}{\draw[coilw] (1.60,{-4.80+0.32*\i}) arc(270:90:0.16);}
\draw[coilw] (1.60,-3.52) -- (1.60,-3.37);
% right coil L2
\draw[coilw] (2.80,-4.95) -- (2.80,-4.80);
\foreach \i in {0,1,2,3}{\draw[coilw] (2.80,{-4.80+0.32*\i}) arc(270:450:0.16);}
\draw[coilw] (2.80,-3.52) -- (2.80,-3.37);
% shared core / coupling flux
\draw[draw=casNedge, dash pattern=on 2pt off 1.6pt, line width=0.6pt]
  (2.14,-5.00) -- (2.14,-3.32);
\draw[draw=casNedge, dash pattern=on 2pt off 1.6pt, line width=0.6pt]
  (2.26,-5.00) -- (2.26,-3.32);

\node[tny] at (1.22,-4.16) {$L_1$};
\node[tny] at (3.18,-4.16) {$L_2$};
\node[tny, fill=white, inner sep=1.5pt] at (2.20,-4.16) {$M$};
\node[tny, anchor=north] at (2.20,-5.12) {$M = k\sqrt{L_1 L_2}$};
\node[tny, anchor=north] at (2.20,-5.50) {$k$ = coupling,\ $0 \le k < 1$};

% =============== (c) flyback diode ==================================
\node[panlbl] at (5.05,-2.94) {(c)\ \ FLYBACK DIODE};

% coil on the right leg
\draw[coilw] (7.60,-4.95) -- (7.60,-4.80);
\foreach \i in {0,1,2,3}{\draw[coilw] (7.60,{-4.80+0.32*\i}) arc(270:450:0.16);}
\draw[coilw] (7.60,-3.52) -- (7.60,-3.37);
\node[tny] at (8.00,-4.16) {$L$};

% loop rails
\draw[wire] (5.95,-3.37) -- (7.60,-3.37);
\draw[wire] (5.95,-4.95) -- (7.60,-4.95);
\fill (7.60,-3.37) circle (1.1pt);
\fill (7.60,-4.95) circle (1.1pt);

% diode: anode low, cathode bar high -- reverse-blocking across the coil
\draw[wire] (5.95,-4.95) -- (5.95,-4.46);
\draw[fill=casEfill, draw=casEedge, line width=0.6pt]
  (5.76,-4.46) -- (6.14,-4.46) -- (5.95,-4.12) -- cycle;
\draw[draw=casEedge, line width=1.0pt] (5.73,-4.12) -- (6.17,-4.12);
\draw[wire] (5.95,-4.12) -- (5.95,-3.37);
\node[tny, anchor=east] at (5.66,-4.29) {D};

% circulating decay loop (clockwise: down the coil, up through the diode)
\draw[-{Latex[length=1.8mm,width=1.4mm]}, line width=0.5pt, draw=casNedge]
  (6.78,-3.78) arc(90:-205:0.38);
\node[tny, anchor=north, align=center] at (6.78,-5.12)
  {blocks while current flows;\\conducts on field collapse};

\end{tikzpicture}
